\icode{E-EME-SI-L3-261121-002}
如图所示，在水平面上\(M\)、\(N\)两点固定着两个电荷量相等的正点电荷，\(A\)、\(B\)为\(MN\)上的两点,\(C\)、\(D\)为\(MN\)中垂线上的两点，且到\(MN\)中点\(O\)的距离相等。下列说法正确的是\paren[]
		
\begin{tikzpicture}[>=stealth,font=\scriptsize]
	\draw[fill=red!50] (-2,0)node[black]{\( +\)} circle(.15) (2,0)node[black]{\( +\)} circle(.15);
	\fill (-1.5,0)node[below]{\(A\)} circle(.05) (1.5,0)node[below]{ \(B\)} circle(.05) (0,1.5)node[right]{ \(C\)} circle(.05) (0,-1.5)node[right]{ \(D\)} circle(.05) (0,0) circle(.05);
	\draw[dashed,cyan] (0,2)--(0,-2) (-1.8,0)--(1.8,0);
	\node at(-.2,-.2){ \(O\)};
	\node at(-2.15,0)[left]{ \(M\)};
	\node at(2.15,0)[right]{ \(N\)};
\end{tikzpicture}
\begin{choices}
	\item  沿中垂线从\(C\)点到\(D\)点电场强度一定先减小后增大
	\item 沿\(MN\)从\(A\)点到\(B\)点电势先升高后降低
	\item 将电子在\(C\)点释放，电子将在\(C\)、\(D\)之间做简谐运动
	\item 将电子沿直线从\(A\)点移到\(O\)点，再从\(O\)点移到\(C\)点，电子的电势能一直增加
\end{choices}
\begin{answer}
	D\qquad 提示：
\end{answer}